%% ============================================================
%%  Acronymes — article « EWC INT8 sur MCU » (version française)
%%  Léonard Rivals — Sprint 40, CL-Embedded
%%
%%  Usage : \usepackage[acronym,toc]{glossaries} puis \makeglossaries
%%          %% ============================================================
%%  Acronymes — article « EWC INT8 sur MCU » (version française)
%%  Léonard Rivals — Sprint 40, CL-Embedded
%%
%%  Usage : \usepackage[acronym,toc]{glossaries} puis \makeglossaries
%%          %% ============================================================
%%  Acronymes — article « EWC INT8 sur MCU » (version française)
%%  Léonard Rivals — Sprint 40, CL-Embedded
%%
%%  Usage : \usepackage[acronym,toc]{glossaries} puis \makeglossaries
%%          %% ============================================================
%%  Acronymes — article « EWC INT8 sur MCU » (version française)
%%  Léonard Rivals — Sprint 40, CL-Embedded
%%
%%  Usage : \usepackage[acronym,toc]{glossaries} puis \makeglossaries
%%          \input{glossary_fr}
%%          \printglossary[type=\acronymtype]
%%
%%  Patron repris du manuscrit (docs/rapport_de_stage/.../glossary_entries.tex),
%%  restreint aux sigles réellement employés dans cet article.
%% ============================================================

%% Les descriptions citent les autres sigles en clair, jamais via \gls{...} :
%% une entrée non citée dans le texte n'est pas imprimée, et le lien pointerait
%% alors vers une ancre inexistante.
\glsnoexpandfields

%% Style : « forme longue (SIGLE) » à la première occurrence, description dans la liste.
\setacronymstyle{long-short-desc}

% ------------------------------------------------------------------
%  APPRENTISSAGE CONTINU
% ------------------------------------------------------------------
\newacronym[description={Apprentissage d'une suite de tâches sans accès aux données
  passées, où le modèle doit intégrer le nouveau sans effacer l'ancien}]%
  {CL}{CL}{Continual Learning}
\newacronym[description={Méthode anti-oubli qui ajoute à la perte une pénalité quadratique
  freinant la modification des poids jugés importants pour les tâches passées, l'importance
  étant estimée par l'information de Fisher. Elle ne stocke aucune donnée passée, ce qui la
  rend compatible avec un microcontrôleur}]%
  {EWC}{EWC}{Elastic Weight Consolidation}
\newacronym[description={Algorithme d'optimisation mettant à jour les poids dans la
  direction opposée au gradient, estimé sur un petit lot d'exemples ; c'est la mise à jour
  exécutée à bord dans le protocole en ligne}]%
  {SGD}{SGD}{Stochastic Gradient Descent}
\newacronym[description={Approche où un réseau préentraîné est gelé en Flash et où seule une
  couche entraînable légère est mise à jour en RAM, échantillon par échantillon}]%
  {TinyOL}{TinyOL}{Tiny Online Learning}
\newacronym[description={Paradigme non neuronal encodant les signaux dans des vecteurs de
  très haute dimension ; l'apprentissage s'y réduit à des opérations élémentaires
  (addition, vote majoritaire), sans rétropropagation}]%
  {HDC}{HDC}{Hyperdimensional Computing}
\newacronym[description={Apprentissage automatique sur microcontrôleur, caractérisé par
  trois contraintes cumulatives : pas de mémoire externe, pas de système d'exploitation,
  et une mémoire vive inférieure à 256\,Ko}]%
  {TinyML}{TinyML}{Tiny Machine Learning}

% ------------------------------------------------------------------
%  QUANTIFICATION
% ------------------------------------------------------------------
\newacronym[description={Quantification appliquée à un modèle déjà entraîné : simple à
  mettre en \oe uvre, mais sujette à des pertes lorsque l'échelle est fixée sans
  calibration ou que la dynamique des tenseurs est large. C'est le régime dont cet article
  mesure l'effondrement puis la récupération}]%
  {PTQ}{PTQ}{Post-Training Quantization}
\newacronym[description={Quantification simulée à chaque pas d'entraînement par
  \emph{fake-quant}, de sorte que le modèle apprenne directement sous contrainte de
  précision réduite ; plus coûteuse à l'entraînement, elle n'apporte ici aucun gain décisif
  face à une PTQ correctement calibrée}]%
  {QAT}{QAT}{Quantization-Aware Training}

% ------------------------------------------------------------------
%  CIBLE MATÉRIELLE ET MÉMOIRE
% ------------------------------------------------------------------
\newacronym[description={Processeur intégrant sur une même puce cœur, mémoire et
  périphériques, destiné aux systèmes embarqués autonomes}]%
  {MCU}{MCU}{Microcontroller Unit}
\newacronym[description={Mémoire vive volatile où résident les variables et les tampons du
  programme ; ressource la plus critique de la cible}]%
  {RAM}{RAM}{Random Access Memory}
\newacronym[description={Mémoire vive statique intégrée au microcontrôleur, rapide mais
  coûteuse en surface de silicium, donc disponible en très faible quantité. C'est elle qui
  fixe le budget mémoire de ce travail}]%
  {SRAM}{SRAM}{Static Random Access Memory}
\newacronym[description={Mémoire rapide directement rattachée au cœur du processeur, non
  accessible à certains périphériques ; sur la carte cible elle complète la mémoire vive
  statique principale}]%
  {CCM}{CCM}{Core Coupled Memory}
\newacronym[description={Circuit spécialisé dans les opérations de réseaux de neurones,
  capable d'exécuter des multiplications-accumulations entières massivement en parallèle.
  Son absence sur la carte cible explique qu'INT8 n'y apporte aucun gain de vitesse}]%
  {NPU}{NPU}{Neural Processing Unit}
\newacronym[description={Unité matérielle de calcul en virgule flottante ; sur la carte
  cible elle rend les opérations FP32 aussi rapides que les opérations entières, ce qui
  annule le bénéfice de vitesse attendu de la quantification}]%
  {FPU}{FPU}{Floating-Point Unit}
\newacronym[description={Jeu d'instructions appliquant une même opération à plusieurs
  données en un cycle ; c'est le levier qui manque au chemin entier de ce travail pour
  convertir le gain de format en gain de temps}]%
  {SIMD}{SIMD}{Single Instruction, Multiple Data}
\newacronym[description={Bibliothèque de noyaux de réseaux de neurones optimisés pour les
  cœurs Cortex-M, exploitant les instructions vectorielles ; piste directe pour lever le
  paradoxe de latence constaté ici}]%
  {CMSISNN}{CMSIS-NN}{Common Microcontroller Software Interface Standard --- Neural Network}
\newacronym[description={Instruction mettant le cœur en sommeil jusqu'à la prochaine
  interruption ; son emploi dans l'attente de la liaison série abaisse le courant de repos
  et rend exploitable la mesure d'énergie par différence}]%
  {WFI}{WFI}{Wait For Interrupt}

% ------------------------------------------------------------------
%  INSTRUMENTATION ET LIAISON
% ------------------------------------------------------------------
\newacronym[description={Unité de débogage du cœur Cortex-M dont le compteur de cycles
  permet de chronométrer une section de code au cycle près, puis de la convertir en
  microsecondes à partir de la fréquence d'horloge. Toutes les latences de cet article en
  sont issues}]%
  {DWT}{DWT}{Data Watchpoint and Trace}
\newacronym[description={Liaison série asynchrone reliant la carte à l'hôte ; elle achemine
  les échantillons trame par trame et rapporte les prédictions}]%
  {UART}{UART}{Universal Asynchronous Receiver-Transmitter}
\newacronym[description={Somme de contrôle calculée sur chaque trame transmise ; un taux
  d'erreur nul garantit que les prédictions comparées sont bien celles produites par la
  carte et non le produit d'une transmission corrompue}]%
  {CRC}{CRC}{Cyclic Redundancy Check}

% ------------------------------------------------------------------
%  COÛT DE CALCUL ET MÉTRIQUES
% ------------------------------------------------------------------
\newacronym[description={Opération élémentaire multipliant deux valeurs et ajoutant le
  produit à un accumulateur ; unité de compte du coût d'une couche dense}]%
  {MAC}{MAC}{Multiply-Accumulate}
\newacronym[description={Opération élémentaire en virgule flottante ; une
  multiplication-accumulation en compte deux}]%
  {FLOP}{FLOP}{Floating-Point Operation}
\newacronym[description={Coût exprimé au niveau du bit, proportionnel au produit des largeurs
  des opérandes ; il prédit un facteur 16 entre FP32 et INT8 que la carte ne restitue pas,
  ce qui illustre les limites d'une métrique analytique}]%
  {BOP}{BOP}{Bit Operation}
\newacronym[description={Aire sous la courbe ROC : probabilité qu'un échantillon anormal
  reçoive un score plus élevé qu'un échantillon normal. Indépendante du seuil de décision}]%
  {AUROC}{AUROC}{Area Under the ROC Curve}

%%          \printglossary[type=\acronymtype]
%%
%%  Patron repris du manuscrit (docs/rapport_de_stage/.../glossary_entries.tex),
%%  restreint aux sigles réellement employés dans cet article.
%% ============================================================

%% Les descriptions citent les autres sigles en clair, jamais via \gls{...} :
%% une entrée non citée dans le texte n'est pas imprimée, et le lien pointerait
%% alors vers une ancre inexistante.
\glsnoexpandfields

%% Style : « forme longue (SIGLE) » à la première occurrence, description dans la liste.
\setacronymstyle{long-short-desc}

% ------------------------------------------------------------------
%  APPRENTISSAGE CONTINU
% ------------------------------------------------------------------
\newacronym[description={Apprentissage d'une suite de tâches sans accès aux données
  passées, où le modèle doit intégrer le nouveau sans effacer l'ancien}]%
  {CL}{CL}{Continual Learning}
\newacronym[description={Méthode anti-oubli qui ajoute à la perte une pénalité quadratique
  freinant la modification des poids jugés importants pour les tâches passées, l'importance
  étant estimée par l'information de Fisher. Elle ne stocke aucune donnée passée, ce qui la
  rend compatible avec un microcontrôleur}]%
  {EWC}{EWC}{Elastic Weight Consolidation}
\newacronym[description={Algorithme d'optimisation mettant à jour les poids dans la
  direction opposée au gradient, estimé sur un petit lot d'exemples ; c'est la mise à jour
  exécutée à bord dans le protocole en ligne}]%
  {SGD}{SGD}{Stochastic Gradient Descent}
\newacronym[description={Approche où un réseau préentraîné est gelé en Flash et où seule une
  couche entraînable légère est mise à jour en RAM, échantillon par échantillon}]%
  {TinyOL}{TinyOL}{Tiny Online Learning}
\newacronym[description={Paradigme non neuronal encodant les signaux dans des vecteurs de
  très haute dimension ; l'apprentissage s'y réduit à des opérations élémentaires
  (addition, vote majoritaire), sans rétropropagation}]%
  {HDC}{HDC}{Hyperdimensional Computing}
\newacronym[description={Apprentissage automatique sur microcontrôleur, caractérisé par
  trois contraintes cumulatives : pas de mémoire externe, pas de système d'exploitation,
  et une mémoire vive inférieure à 256\,Ko}]%
  {TinyML}{TinyML}{Tiny Machine Learning}

% ------------------------------------------------------------------
%  QUANTIFICATION
% ------------------------------------------------------------------
\newacronym[description={Quantification appliquée à un modèle déjà entraîné : simple à
  mettre en \oe uvre, mais sujette à des pertes lorsque l'échelle est fixée sans
  calibration ou que la dynamique des tenseurs est large. C'est le régime dont cet article
  mesure l'effondrement puis la récupération}]%
  {PTQ}{PTQ}{Post-Training Quantization}
\newacronym[description={Quantification simulée à chaque pas d'entraînement par
  \emph{fake-quant}, de sorte que le modèle apprenne directement sous contrainte de
  précision réduite ; plus coûteuse à l'entraînement, elle n'apporte ici aucun gain décisif
  face à une PTQ correctement calibrée}]%
  {QAT}{QAT}{Quantization-Aware Training}

% ------------------------------------------------------------------
%  CIBLE MATÉRIELLE ET MÉMOIRE
% ------------------------------------------------------------------
\newacronym[description={Processeur intégrant sur une même puce cœur, mémoire et
  périphériques, destiné aux systèmes embarqués autonomes}]%
  {MCU}{MCU}{Microcontroller Unit}
\newacronym[description={Mémoire vive volatile où résident les variables et les tampons du
  programme ; ressource la plus critique de la cible}]%
  {RAM}{RAM}{Random Access Memory}
\newacronym[description={Mémoire vive statique intégrée au microcontrôleur, rapide mais
  coûteuse en surface de silicium, donc disponible en très faible quantité. C'est elle qui
  fixe le budget mémoire de ce travail}]%
  {SRAM}{SRAM}{Static Random Access Memory}
\newacronym[description={Mémoire rapide directement rattachée au cœur du processeur, non
  accessible à certains périphériques ; sur la carte cible elle complète la mémoire vive
  statique principale}]%
  {CCM}{CCM}{Core Coupled Memory}
\newacronym[description={Circuit spécialisé dans les opérations de réseaux de neurones,
  capable d'exécuter des multiplications-accumulations entières massivement en parallèle.
  Son absence sur la carte cible explique qu'INT8 n'y apporte aucun gain de vitesse}]%
  {NPU}{NPU}{Neural Processing Unit}
\newacronym[description={Unité matérielle de calcul en virgule flottante ; sur la carte
  cible elle rend les opérations FP32 aussi rapides que les opérations entières, ce qui
  annule le bénéfice de vitesse attendu de la quantification}]%
  {FPU}{FPU}{Floating-Point Unit}
\newacronym[description={Jeu d'instructions appliquant une même opération à plusieurs
  données en un cycle ; c'est le levier qui manque au chemin entier de ce travail pour
  convertir le gain de format en gain de temps}]%
  {SIMD}{SIMD}{Single Instruction, Multiple Data}
\newacronym[description={Bibliothèque de noyaux de réseaux de neurones optimisés pour les
  cœurs Cortex-M, exploitant les instructions vectorielles ; piste directe pour lever le
  paradoxe de latence constaté ici}]%
  {CMSISNN}{CMSIS-NN}{Common Microcontroller Software Interface Standard --- Neural Network}
\newacronym[description={Instruction mettant le cœur en sommeil jusqu'à la prochaine
  interruption ; son emploi dans l'attente de la liaison série abaisse le courant de repos
  et rend exploitable la mesure d'énergie par différence}]%
  {WFI}{WFI}{Wait For Interrupt}

% ------------------------------------------------------------------
%  INSTRUMENTATION ET LIAISON
% ------------------------------------------------------------------
\newacronym[description={Unité de débogage du cœur Cortex-M dont le compteur de cycles
  permet de chronométrer une section de code au cycle près, puis de la convertir en
  microsecondes à partir de la fréquence d'horloge. Toutes les latences de cet article en
  sont issues}]%
  {DWT}{DWT}{Data Watchpoint and Trace}
\newacronym[description={Liaison série asynchrone reliant la carte à l'hôte ; elle achemine
  les échantillons trame par trame et rapporte les prédictions}]%
  {UART}{UART}{Universal Asynchronous Receiver-Transmitter}
\newacronym[description={Somme de contrôle calculée sur chaque trame transmise ; un taux
  d'erreur nul garantit que les prédictions comparées sont bien celles produites par la
  carte et non le produit d'une transmission corrompue}]%
  {CRC}{CRC}{Cyclic Redundancy Check}

% ------------------------------------------------------------------
%  COÛT DE CALCUL ET MÉTRIQUES
% ------------------------------------------------------------------
\newacronym[description={Opération élémentaire multipliant deux valeurs et ajoutant le
  produit à un accumulateur ; unité de compte du coût d'une couche dense}]%
  {MAC}{MAC}{Multiply-Accumulate}
\newacronym[description={Opération élémentaire en virgule flottante ; une
  multiplication-accumulation en compte deux}]%
  {FLOP}{FLOP}{Floating-Point Operation}
\newacronym[description={Coût exprimé au niveau du bit, proportionnel au produit des largeurs
  des opérandes ; il prédit un facteur 16 entre FP32 et INT8 que la carte ne restitue pas,
  ce qui illustre les limites d'une métrique analytique}]%
  {BOP}{BOP}{Bit Operation}
\newacronym[description={Aire sous la courbe ROC : probabilité qu'un échantillon anormal
  reçoive un score plus élevé qu'un échantillon normal. Indépendante du seuil de décision}]%
  {AUROC}{AUROC}{Area Under the ROC Curve}

%%          \printglossary[type=\acronymtype]
%%
%%  Patron repris du manuscrit (docs/rapport_de_stage/.../glossary_entries.tex),
%%  restreint aux sigles réellement employés dans cet article.
%% ============================================================

%% Les descriptions citent les autres sigles en clair, jamais via \gls{...} :
%% une entrée non citée dans le texte n'est pas imprimée, et le lien pointerait
%% alors vers une ancre inexistante.
\glsnoexpandfields

%% Style : « forme longue (SIGLE) » à la première occurrence, description dans la liste.
\setacronymstyle{long-short-desc}

% ------------------------------------------------------------------
%  APPRENTISSAGE CONTINU
% ------------------------------------------------------------------
\newacronym[description={Apprentissage d'une suite de tâches sans accès aux données
  passées, où le modèle doit intégrer le nouveau sans effacer l'ancien}]%
  {CL}{CL}{Continual Learning}
\newacronym[description={Méthode anti-oubli qui ajoute à la perte une pénalité quadratique
  freinant la modification des poids jugés importants pour les tâches passées, l'importance
  étant estimée par l'information de Fisher. Elle ne stocke aucune donnée passée, ce qui la
  rend compatible avec un microcontrôleur}]%
  {EWC}{EWC}{Elastic Weight Consolidation}
\newacronym[description={Algorithme d'optimisation mettant à jour les poids dans la
  direction opposée au gradient, estimé sur un petit lot d'exemples ; c'est la mise à jour
  exécutée à bord dans le protocole en ligne}]%
  {SGD}{SGD}{Stochastic Gradient Descent}
\newacronym[description={Approche où un réseau préentraîné est gelé en Flash et où seule une
  couche entraînable légère est mise à jour en RAM, échantillon par échantillon}]%
  {TinyOL}{TinyOL}{Tiny Online Learning}
\newacronym[description={Paradigme non neuronal encodant les signaux dans des vecteurs de
  très haute dimension ; l'apprentissage s'y réduit à des opérations élémentaires
  (addition, vote majoritaire), sans rétropropagation}]%
  {HDC}{HDC}{Hyperdimensional Computing}
\newacronym[description={Apprentissage automatique sur microcontrôleur, caractérisé par
  trois contraintes cumulatives : pas de mémoire externe, pas de système d'exploitation,
  et une mémoire vive inférieure à 256\,Ko}]%
  {TinyML}{TinyML}{Tiny Machine Learning}

% ------------------------------------------------------------------
%  QUANTIFICATION
% ------------------------------------------------------------------
\newacronym[description={Quantification appliquée à un modèle déjà entraîné : simple à
  mettre en \oe uvre, mais sujette à des pertes lorsque l'échelle est fixée sans
  calibration ou que la dynamique des tenseurs est large. C'est le régime dont cet article
  mesure l'effondrement puis la récupération}]%
  {PTQ}{PTQ}{Post-Training Quantization}
\newacronym[description={Quantification simulée à chaque pas d'entraînement par
  \emph{fake-quant}, de sorte que le modèle apprenne directement sous contrainte de
  précision réduite ; plus coûteuse à l'entraînement, elle n'apporte ici aucun gain décisif
  face à une PTQ correctement calibrée}]%
  {QAT}{QAT}{Quantization-Aware Training}

% ------------------------------------------------------------------
%  CIBLE MATÉRIELLE ET MÉMOIRE
% ------------------------------------------------------------------
\newacronym[description={Processeur intégrant sur une même puce cœur, mémoire et
  périphériques, destiné aux systèmes embarqués autonomes}]%
  {MCU}{MCU}{Microcontroller Unit}
\newacronym[description={Mémoire vive volatile où résident les variables et les tampons du
  programme ; ressource la plus critique de la cible}]%
  {RAM}{RAM}{Random Access Memory}
\newacronym[description={Mémoire vive statique intégrée au microcontrôleur, rapide mais
  coûteuse en surface de silicium, donc disponible en très faible quantité. C'est elle qui
  fixe le budget mémoire de ce travail}]%
  {SRAM}{SRAM}{Static Random Access Memory}
\newacronym[description={Mémoire rapide directement rattachée au cœur du processeur, non
  accessible à certains périphériques ; sur la carte cible elle complète la mémoire vive
  statique principale}]%
  {CCM}{CCM}{Core Coupled Memory}
\newacronym[description={Circuit spécialisé dans les opérations de réseaux de neurones,
  capable d'exécuter des multiplications-accumulations entières massivement en parallèle.
  Son absence sur la carte cible explique qu'INT8 n'y apporte aucun gain de vitesse}]%
  {NPU}{NPU}{Neural Processing Unit}
\newacronym[description={Unité matérielle de calcul en virgule flottante ; sur la carte
  cible elle rend les opérations FP32 aussi rapides que les opérations entières, ce qui
  annule le bénéfice de vitesse attendu de la quantification}]%
  {FPU}{FPU}{Floating-Point Unit}
\newacronym[description={Jeu d'instructions appliquant une même opération à plusieurs
  données en un cycle ; c'est le levier qui manque au chemin entier de ce travail pour
  convertir le gain de format en gain de temps}]%
  {SIMD}{SIMD}{Single Instruction, Multiple Data}
\newacronym[description={Bibliothèque de noyaux de réseaux de neurones optimisés pour les
  cœurs Cortex-M, exploitant les instructions vectorielles ; piste directe pour lever le
  paradoxe de latence constaté ici}]%
  {CMSISNN}{CMSIS-NN}{Common Microcontroller Software Interface Standard --- Neural Network}
\newacronym[description={Instruction mettant le cœur en sommeil jusqu'à la prochaine
  interruption ; son emploi dans l'attente de la liaison série abaisse le courant de repos
  et rend exploitable la mesure d'énergie par différence}]%
  {WFI}{WFI}{Wait For Interrupt}

% ------------------------------------------------------------------
%  INSTRUMENTATION ET LIAISON
% ------------------------------------------------------------------
\newacronym[description={Unité de débogage du cœur Cortex-M dont le compteur de cycles
  permet de chronométrer une section de code au cycle près, puis de la convertir en
  microsecondes à partir de la fréquence d'horloge. Toutes les latences de cet article en
  sont issues}]%
  {DWT}{DWT}{Data Watchpoint and Trace}
\newacronym[description={Liaison série asynchrone reliant la carte à l'hôte ; elle achemine
  les échantillons trame par trame et rapporte les prédictions}]%
  {UART}{UART}{Universal Asynchronous Receiver-Transmitter}
\newacronym[description={Somme de contrôle calculée sur chaque trame transmise ; un taux
  d'erreur nul garantit que les prédictions comparées sont bien celles produites par la
  carte et non le produit d'une transmission corrompue}]%
  {CRC}{CRC}{Cyclic Redundancy Check}

% ------------------------------------------------------------------
%  COÛT DE CALCUL ET MÉTRIQUES
% ------------------------------------------------------------------
\newacronym[description={Opération élémentaire multipliant deux valeurs et ajoutant le
  produit à un accumulateur ; unité de compte du coût d'une couche dense}]%
  {MAC}{MAC}{Multiply-Accumulate}
\newacronym[description={Opération élémentaire en virgule flottante ; une
  multiplication-accumulation en compte deux}]%
  {FLOP}{FLOP}{Floating-Point Operation}
\newacronym[description={Coût exprimé au niveau du bit, proportionnel au produit des largeurs
  des opérandes ; il prédit un facteur 16 entre FP32 et INT8 que la carte ne restitue pas,
  ce qui illustre les limites d'une métrique analytique}]%
  {BOP}{BOP}{Bit Operation}
\newacronym[description={Aire sous la courbe ROC : probabilité qu'un échantillon anormal
  reçoive un score plus élevé qu'un échantillon normal. Indépendante du seuil de décision}]%
  {AUROC}{AUROC}{Area Under the ROC Curve}

%%          \printglossary[type=\acronymtype]
%%
%%  Patron repris du manuscrit (docs/rapport_de_stage/.../glossary_entries.tex),
%%  restreint aux sigles réellement employés dans cet article.
%% ============================================================

%% Les descriptions citent les autres sigles en clair, jamais via \gls{...} :
%% une entrée non citée dans le texte n'est pas imprimée, et le lien pointerait
%% alors vers une ancre inexistante.
\glsnoexpandfields

%% Style : « forme longue (SIGLE) » à la première occurrence, description dans la liste.
\setacronymstyle{long-short-desc}

% ------------------------------------------------------------------
%  APPRENTISSAGE CONTINU
% ------------------------------------------------------------------
\newacronym[description={Apprentissage d'une suite de tâches sans accès aux données
  passées, où le modèle doit intégrer le nouveau sans effacer l'ancien}]%
  {CL}{CL}{Continual Learning}
\newacronym[description={Méthode anti-oubli qui ajoute à la perte une pénalité quadratique
  freinant la modification des poids jugés importants pour les tâches passées, l'importance
  étant estimée par l'information de Fisher. Elle ne stocke aucune donnée passée, ce qui la
  rend compatible avec un microcontrôleur}]%
  {EWC}{EWC}{Elastic Weight Consolidation}
\newacronym[description={Algorithme d'optimisation mettant à jour les poids dans la
  direction opposée au gradient, estimé sur un petit lot d'exemples ; c'est la mise à jour
  exécutée à bord dans le protocole en ligne}]%
  {SGD}{SGD}{Stochastic Gradient Descent}
\newacronym[description={Approche où un réseau préentraîné est gelé en Flash et où seule une
  couche entraînable légère est mise à jour en RAM, échantillon par échantillon}]%
  {TinyOL}{TinyOL}{Tiny Online Learning}
\newacronym[description={Paradigme non neuronal encodant les signaux dans des vecteurs de
  très haute dimension ; l'apprentissage s'y réduit à des opérations élémentaires
  (addition, vote majoritaire), sans rétropropagation}]%
  {HDC}{HDC}{Hyperdimensional Computing}
\newacronym[description={Apprentissage automatique sur microcontrôleur, caractérisé par
  trois contraintes cumulatives : pas de mémoire externe, pas de système d'exploitation,
  et une mémoire vive inférieure à 256\,Ko}]%
  {TinyML}{TinyML}{Tiny Machine Learning}

% ------------------------------------------------------------------
%  QUANTIFICATION
% ------------------------------------------------------------------
\newacronym[description={Quantification appliquée à un modèle déjà entraîné : simple à
  mettre en \oe uvre, mais sujette à des pertes lorsque l'échelle est fixée sans
  calibration ou que la dynamique des tenseurs est large. C'est le régime dont cet article
  mesure l'effondrement puis la récupération}]%
  {PTQ}{PTQ}{Post-Training Quantization}
\newacronym[description={Quantification simulée à chaque pas d'entraînement par
  \emph{fake-quant}, de sorte que le modèle apprenne directement sous contrainte de
  précision réduite ; plus coûteuse à l'entraînement, elle n'apporte ici aucun gain décisif
  face à une PTQ correctement calibrée}]%
  {QAT}{QAT}{Quantization-Aware Training}

% ------------------------------------------------------------------
%  CIBLE MATÉRIELLE ET MÉMOIRE
% ------------------------------------------------------------------
\newacronym[description={Processeur intégrant sur une même puce cœur, mémoire et
  périphériques, destiné aux systèmes embarqués autonomes}]%
  {MCU}{MCU}{Microcontroller Unit}
\newacronym[description={Mémoire vive volatile où résident les variables et les tampons du
  programme ; ressource la plus critique de la cible}]%
  {RAM}{RAM}{Random Access Memory}
\newacronym[description={Mémoire vive statique intégrée au microcontrôleur, rapide mais
  coûteuse en surface de silicium, donc disponible en très faible quantité. C'est elle qui
  fixe le budget mémoire de ce travail}]%
  {SRAM}{SRAM}{Static Random Access Memory}
\newacronym[description={Mémoire rapide directement rattachée au cœur du processeur, non
  accessible à certains périphériques ; sur la carte cible elle complète la mémoire vive
  statique principale}]%
  {CCM}{CCM}{Core Coupled Memory}
\newacronym[description={Circuit spécialisé dans les opérations de réseaux de neurones,
  capable d'exécuter des multiplications-accumulations entières massivement en parallèle.
  Son absence sur la carte cible explique qu'INT8 n'y apporte aucun gain de vitesse}]%
  {NPU}{NPU}{Neural Processing Unit}
\newacronym[description={Unité matérielle de calcul en virgule flottante ; sur la carte
  cible elle rend les opérations FP32 aussi rapides que les opérations entières, ce qui
  annule le bénéfice de vitesse attendu de la quantification}]%
  {FPU}{FPU}{Floating-Point Unit}
\newacronym[description={Jeu d'instructions appliquant une même opération à plusieurs
  données en un cycle ; c'est le levier qui manque au chemin entier de ce travail pour
  convertir le gain de format en gain de temps}]%
  {SIMD}{SIMD}{Single Instruction, Multiple Data}
\newacronym[description={Bibliothèque de noyaux de réseaux de neurones optimisés pour les
  cœurs Cortex-M, exploitant les instructions vectorielles ; piste directe pour lever le
  paradoxe de latence constaté ici}]%
  {CMSISNN}{CMSIS-NN}{Common Microcontroller Software Interface Standard --- Neural Network}
\newacronym[description={Instruction mettant le cœur en sommeil jusqu'à la prochaine
  interruption ; son emploi dans l'attente de la liaison série abaisse le courant de repos
  et rend exploitable la mesure d'énergie par différence}]%
  {WFI}{WFI}{Wait For Interrupt}

% ------------------------------------------------------------------
%  INSTRUMENTATION ET LIAISON
% ------------------------------------------------------------------
\newacronym[description={Unité de débogage du cœur Cortex-M dont le compteur de cycles
  permet de chronométrer une section de code au cycle près, puis de la convertir en
  microsecondes à partir de la fréquence d'horloge. Toutes les latences de cet article en
  sont issues}]%
  {DWT}{DWT}{Data Watchpoint and Trace}
\newacronym[description={Liaison série asynchrone reliant la carte à l'hôte ; elle achemine
  les échantillons trame par trame et rapporte les prédictions}]%
  {UART}{UART}{Universal Asynchronous Receiver-Transmitter}
\newacronym[description={Somme de contrôle calculée sur chaque trame transmise ; un taux
  d'erreur nul garantit que les prédictions comparées sont bien celles produites par la
  carte et non le produit d'une transmission corrompue}]%
  {CRC}{CRC}{Cyclic Redundancy Check}

% ------------------------------------------------------------------
%  COÛT DE CALCUL ET MÉTRIQUES
% ------------------------------------------------------------------
\newacronym[description={Opération élémentaire multipliant deux valeurs et ajoutant le
  produit à un accumulateur ; unité de compte du coût d'une couche dense}]%
  {MAC}{MAC}{Multiply-Accumulate}
\newacronym[description={Opération élémentaire en virgule flottante ; une
  multiplication-accumulation en compte deux}]%
  {FLOP}{FLOP}{Floating-Point Operation}
\newacronym[description={Coût exprimé au niveau du bit, proportionnel au produit des largeurs
  des opérandes ; il prédit un facteur 16 entre FP32 et INT8 que la carte ne restitue pas,
  ce qui illustre les limites d'une métrique analytique}]%
  {BOP}{BOP}{Bit Operation}
\newacronym[description={Aire sous la courbe ROC : probabilité qu'un échantillon anormal
  reçoive un score plus élevé qu'un échantillon normal. Indépendante du seuil de décision}]%
  {AUROC}{AUROC}{Area Under the ROC Curve}
